\documentclass[letterpaper,notitlepage,12pt]{book}
\author{Leo Przybylski}
\usepackage[usenames,dvipsnames]{color}
\usepackage{minted}
\usepackage{graphicx}
\usepackage{epstopdf}
\usepackage{listings}
\usepackage{fancyhdr}
\usepackage{hyperref}
\hypersetup{
    colorlinks,
    citecolor=ForestGreen,
    filecolor=ForestGreen,
    linkcolor=ForestGreen,
    urlcolor=ForestGreen
}
\hypersetup{linktocpage}
%%\definecolor{ubergray}{RGB}{245,245,245}
\title{\includegraphics[width=\textwidth]{cover.png}\\KFS Hands-On Developer Introduction}
\date{}
\pagestyle{fancy}
\fancyhead{} % Clear all header fields 
\fancyhead[OL]{\sectionmark}
\fancyhead[OR]{\includegraphics[height=26pt]{heading.png}}% 
\fancyhead[ER]{\sectionmark}
\fancyhead[EL]{\includegraphics[height=26pt]{heading.png}}% 

\begin{document}
\maketitle
\tableofcontents

\addcontentsline{toc}{part}{Preface}
\part*{Preface}

\addcontentsline{toc}{section}{Copyright}
\section*{Copyright}
\addcontentsline{toc}{subsection}{Copyright Holder}
\subsection*{Copyright Holder}
\copyright Copyright 2011, 2012
Leo Przybylski
leo@rsmart.com

\addcontentsline{toc}{subsection}{Disclaimer}
\subsection*{Disclaimer}
This work is licensed under the Creative Commons Attribution-ShareAlike 3.0 United States License. To view a copy of this license, visit http://creativecommons.org/licenses/by-sa/3.0/us/ or send a letter to Creative Commons, 444 Castro Street, Suite 900, Mountain View, California, 94041, USA.

\addcontentsline{toc}{section}{About the Trainer}
\section*{About the Trainer}
Leo started working with the Kuali Foundation in 2005 as a developer
on the Kuali Financial System. Since then, he has worked as a
\emph{Development Manager} on the Kuali Financial System, \emph{Lead Developer} on
the Kuali Coeus project, \emph{Software Architect} on the University of
Arizona KFS implementation, and now is a \emph{Release Engineer} for the
Kuali Foundation for the Rice Project. 

Leo has given six presentations on KFS, KC, and Rice on to separate
Kuali Days occassions. 

One significant contribution he has made to the Kuali Community is his
Rice LDAP Integration module.

\addcontentsline{toc}{section}{Using these Exercises}
\section*{Using these Exercises}
\addcontentsline{toc}{subsection}{VirtualBox Appliance}
\subsection*{VirtualBox Appliance}
Exercise instructions are included in this document. All software and
examples are available on the VirtualBox appliance distributed during
class. To install the VirtualBox appliance:
\begin{enumerate}
  \item Execute the VirtualBox installer to install the software.
  \item Copy the \textbf{KFSDev.box} from the distributed USB drive to
    your hard disk.
  \item Execute from a shell 
    \begin{minted}{sh}
      vagrant box add KFSDev KFSDev.box
    \end{minted}
  \item Run init
    \begin{minted}{sh}
      vagrant init KFSDev
    \end{minted}
  \item Import the Virtual Machine
    \begin{minted}{sh}
      vagrant up
    \end{minted}
  \item Stop the Virtual Machine
    \begin{minted}{sh}
      vagrant halt
    \end{minted}
  \item Start VirtualBox
\end{enumerate}

\addcontentsline{toc}{subsection}{Virtual Machine Manifest}
\subsection*{Virtual Machine Manifest}
The VirtualBox appliance is an Ubuntu Linux distribution. Within it is
the software we will use for this class:
\begin{description}
  \item [Eclipse Indigo] the IDE used for class. Includes Subclipse,
    the m2eclipse plugin, and pre-installed projects with examples.
  \item [OpenJDK 1.7.0\_06 IcedTea] the JVM used for executing/testing
    examples.
  \item [Maven 3] used to build Rice applications, run tests, and
    start the Tomcat6 application
\end{description}

\subsubsection*{Credentials}
\begin{description}
\item [User Account] is \textbf{kuali} with the password
  \textbf{kuali}. This is used to unlock the VM after it has suspended,
  gone to sleep, or locked. The password is also required for
  executing commands as \textbf{root} which may on occassion be
  required. The user account home directory is located at
  \textbf{/home/kuali} and will frequently be referred to during the training.
\item [Database Account] uses the jdbc connection string
  \textbf{jdbc:mysql://localhost:3306/kuldemo} and the
  username/password \textbf{kuldemo}/\textbf{kuldemo}. These are the
  default credentials and database connection information as defined
  in kul-cfg-dbs.
\end{description}

\subsubsection*{Structure}
The Eclipse workspace is located at \textbf{/home/kuali/workspace}. 


\addcontentsline{toc}{section}{Training Overview}
\section*{Training Overview}

\addcontentsline{toc}{part}{Day 3: Advanced Kuali Identity Managemet}
\part*{Day 3\\
Advanced Kuali Identity Management}

\addcontentsline{toc}{section}{Exercise4: ldapsearch Tool and Query Syntax}
{\setlength{\baselineskip}%
  {0.0\baselineskip}
  \section*{\flushright Exercise 4\\
ldapsearch Tool and the Query Syntax}
  \hrulefill \par}

\addcontentsline{toc}{subsection}{Description}
\subsection*{Description}

\addcontentsline{toc}{subsection}{Goals}
\subsection*{Goals}

\addcontentsline{toc}{subsection}{Instructions}
\subsection*{Instructions}
\begin{enumerate}
\item Start by verifying the tools are installed by running ldapsearch
  on the commandline by with '-h' to access the help information
\begin{lstlisting}[language=bash,basicstyle=\scriptsize,backgroundcolor=\color{ubergray},caption={Verify
  the ldapsearch is there},frame=single,breaklines=true]
ldapsearch -h
man ldapsearch
\end{lstlisting}
\item Search by connecting to the localhost LDAP server to the
  \textbf{cn-admin,dc=rsmart,dc=com} distinguished name and the
  \textbf{ou=people,dc=rsmart,dc=com} organizational unit, and query
  for anyone with the \textbf{eduPersonPrincipalName} of leo
\begin{lstlisting}[language=bash,basicstyle=\scriptsize,backgroundcolor=\color{ubergray},caption={Verify
  the ldapsearch is there},frame=single,breaklines=true]
ldapsearch -H ldap://localhost -D "cn=admin,dc=rsmart,dc=com" -b
"ou=people,dc=rsmart,dc=com" -w rice
"(eduPersonPrincipalName=leo)"\end{lstlisting}
This query should result in the following:
\begin{lstlisting}[language=bash,basicstyle=\scriptsize,backgroundcolor=\color{ubergray},caption={Basic
  query using eduPersonPrincipalName},frame=single,breaklines=true]
# extended LDIF
#
# LDAPv3
# base <ou=people,dc=rsmart,dc=com> with scope subtree
# filter: (eduPersonPrincipalName=leo)
# requesting: ALL
#

# 10000000, people, rsmart.com
dn: uid=10000000,ou=people,dc=rsmart,dc=com
uid: 10000000
eduPersonPrincipalName: leo
mail: leo@rsmart.com
eduPersonPrimaryAffiliation: student
eduPersonAffiliation: former-employee
eduPersonAffiliation: former-staff
eduPersonAffiliation: member
eduPersonAffiliation: student
departmentNumber: 9507
sn: Przybylski
givenName: Leonard
cn: Leonard Przybylski
employeeStatus: T
employeeZip: 85721-0073
employeeState: AZ
employeeCity: TUCSON
employeePhone: 5206266997
employeeNumber: 133006641
objectClass: top
objectClass: eduPerson
objectClass: organizationalPerson
objectClass: inetOrgPerson
objectClass: rsmartEmployee

# 10000000, people, rsmart.com
dn: cn=10000000,ou=people,dc=rsmart,dc=com
uid: 10000000
eduPersonPrincipalName: leo
mail: leo@rsmart.com
eduPersonPrimaryAffiliation: student
eduPersonAffiliation: former-employee
eduPersonAffiliation: former-staff
eduPersonAffiliation: member
eduPersonAffiliation: student
departmentNumber: 9507
sn: Przybylski
givenName: Leonard
cn: Leonard Przybylski
cn: 10000000
employeeStatus: T
employeeZip: 85721-0073
employeeState: AZ
employeeCity: TUCSON
employeePhone: 5206266997
employeeNumber: 133006641
objectClass: top
objectClass: eduPerson
objectClass: organizationalPerson
objectClass: inetOrgPerson
objectClass: rsmartEmployee

# search result
search: 2
result: 0 Success

# numResponses: 3
\end{lstlisting}

Information you want to analyze that will be used later is that we are
using 
\begin{lstlisting}[language=bash,basicstyle=\scriptsize,backgroundcolor=\color{ubergray},caption={Verify
  the ldapsearch is there},frame=single,breaklines=true]
uid: 10000000
eduPersonPrincipalName: leo
\end{lstlisting}

Later, we will need to adjust our mappings in our integration to be
aware of these fields. Knowing what they are is important. 

Also, notice that we are using the following LDAP objectClasses.
\begin{description}
\item [inetOrgPerson] or Internet Organization Person. This schema
  is a standard in Directory Services containing a number of important
  attribute fields we use in our data like \textbf{employeeNumber},
  \textbf{departmentNumber}, and \textbf{employeeType}.
\item [eduPerson] is an internet2 standard for information about
  higher education people to be shared across institutions. Fields
  from this that we're using are anything starting with \textbf{eduPerson*}
\item [rsmartEmployee] my own schema I created to fill in gaps for
  employees like the \textbf{employeeStatus} field which is actually
  really important.
\end{description}

\item Execute ldapsearch again. This time search for a \textbf{cn} of
  10000000 \textbf{OR} \textbf{mail} field of \verb|*b*@rsmart.com|

\begin{lstlisting}[language=bash,basicstyle=\scriptsize,backgroundcolor=\color{ubergray},caption={LDAP
  Query syntax using OR notation},frame=single,breaklines=true]
ldapsearch -H ldap://localhost -D "cn=admin,dc=rsmart,dc=com" -b
"ou=people,dc=rsmart,dc=com" -w rice
"(|(cn=10000000)(mail=*b*@rsmart.com))"
\end{lstlisting}
As in most query syntaxes, \textbf{OR} and \textbf{AND} are pretty much
universal. What sets LDAP apart is how most criteria is contained
within \textbf{()}. Further, operators are prefixed. That is, they
come before the criteria. Notice how the $|$ comes first in 
\verb/"(|(cn=10000000)(mail=*b*@rsmart.com))"/
\end{enumerate}


\addcontentsline{toc}{section}{Exercise 5: Implementing LDAP Entity Integration}
{\setlength{\baselineskip}%
  {0.0\baselineskip}
  \section*{\flushright Exercise 5\\
  Implementing LDAP Entity Integration}
  \hrulefill \par}

The Kuali Foundation official documentation for the LDAP Integration
Module is located at https://wiki.kuali.org/x/FSyREg

\addcontentsline{toc}{subsection}{Setup CAS for LDAP}
\subsection*{1 Setup CAS for LDAP}
Up until now, we have been using the \verb|DummyLoginFilter| which is
ok, but now that our users exist in CAS, we can no longer use it. Once
CAS validates a users, KIM will try to find the user in the system. We
will need to integrate both steps at once before the application will
function again.

I have included in the workspace of the VM a \textbf{CAS} project. You
will want to import this into Eclipse as a new Maven project.

\begin{enumerate}
\item File $\rightarrow$ Import $\rightarrow$ Other

%%\includegraphics[width=\textwidth]{images/Screenshot-Import.png}

\item Click the \textbf{Browse} button
\item Select \textbf{cas} from \verb|/home/rice/workspace|

%%\includegraphics[width=\textwidth]{images/Screenshot14.png}
\item Edit \verb|src/main/resources/custom.properties|.

%%\includegraphics[width=\textwidth]{images/Screenshot15.png}
 You want it to look like this.
\begin{lstlisting}[language=bash,basicstyle=\scriptsize,backgroundcolor=\color{ubergray},caption={src/main/resources/custom.properties},frame=single,breaklines=true]
ldap.server.url=ldap://localhost
ldap.server.bind.username=cn=admin,dc=rsmart,dc=com
ldap.server.bind.password=rice
ldap.authentication.filter=eduPersonPrincipalName=%u,ou=people,dc=rsmart,dc=com
ldap.searchBase=ou=people,dc=rsmart,dc=com
\end{lstlisting}
\item Start it up with Maven
\begin{lstlisting}[language=bash,basicstyle=\scriptsize,backgroundcolor=\color{ubergray},caption={src/main/resources/custom.properties},frame=single,breaklines=true]
mvn war:inplace tomcat:run
\end{lstlisting}
\item Configure the application for CAS by editing the
  \verb|/home/rice/kuali/main/dev/trnapp-config.xml|
\begin{lstlisting}[numbers=left,language=xml,basicstyle=\scriptsize,backgroundcolor=\color{ubergray},caption={src/main/resources/custom.properties},frame=single,breaklines=true]
	<param name="rice.ldap.username">cn=admin,dc=rsmart,dc=com</param>
    <param name="rice.ldap.password">rice</param>
    <param name="rice.ldap.url">ldap://localhost</param>
    <param name="rice.ldap.base">ou=people,dc=rsmart,dc=com</param>
	<param name="rice.additionalSpringFiles">org/kuali/rice/kim/config/KIMLdapSpringBeans.xml</param>

      <param name="cas.rice.server.name">${application.host}:8081</param>
      <param name="cas.url">http://${cas.rice.server.name}/${cas.context.name}</param>
      <param name="cas.require.https">false</param>
      <param name="cas.validate.password">false</param>
       <param name="filter.login.class">org.jasig.cas.client.authentication.AuthenticationFilter</param>
       <param name="filter.login.casServerLoginUrl">${cas.url}/login</param>
       <param name="filter.login.serverName">${application.host}:${http.port}</param>
       <param name="filtermapping.login.1">/*</param>

       <param name="filter.validation.class">org.jasig.cas.client.validation.Cas20ProxyReceivingTicketValidationFilter</param>
       <param name="filter.validation.casServerUrlPrefix">${cas.url}</param>
       <param name="filter.validation.serverName">${application.host}:${http.port}</param>
       <param name="filtermapping.validation.2">/*</param>

       <param name="filter.caswrapper.class">org.jasig.cas.client.util.HttpServletRequestWrapperFilter</param>
       <param name="filtermapping.caswrapper.3">/*</param>
\end{lstlisting}
\item Browse to http://localhost:8081/cas/login

%%\includegraphics[width=\textwidth]{images/Screenshot16.png}

\end{enumerate}

\addcontentsline{toc}{subsection}{Add the Rice LDAP Integration Module}
\subsection*{2 Add the Rice LDAP Integration Module}
Add the LDAP Module to a Rice project and customize the mappers.

\addcontentsline{toc}{subsection}{Goals}
\subsection*{Goals}
\begin{enumerate}
  \item Attache the LDAP module to Rice 1.0.3 and build it
  \item Learn to connect CAS to LDAP
  \item Learn to debug and troubleshoot LDAP Queries
  \item Learn to map between an LDAP database and KIM objects
  \item Learn to modify mappings at all levels
  \item Learn about caching in KIM
  \item Learn about configuring Rice LDAP Integration
\end{enumerate}

\addcontentsline{toc}{subsection}{Implementation}
\subsection*{1 Implementation}
\subsubsection*{1.1 Enable LDAP Integration}
\begin{enumerate}
\item Switch to the \verb|rice-src/ldap/| project
\item Browse to \verb|src/main/config/example-config/|
\item Open the rice-config.xml
\item Copy the contents into \verb|trnapp-config.xml|
\item Change the fields until they are like this:
\begin{lstlisting}[numbers=left,language=xml,basicstyle=\scriptsize,backgroundcolor=\color{ubergray},caption={trnapp-config.xml},frame=single,breaklines=true]
  <param name="rice.ldap.username">cn=admin,ou=people,dc=rsmart,dc=com</param>
  <param name="rice.ldap.password">rice</param>
  <param name="rice.ldap.url">ldap://localhost</param>
  <param name="rice.ldap.base">ou=People,dc=rsmart,dc=edu</param>
  <param
  name="rice.additionalSpringFiles">org/kuali/rice/kim/config/KIMLdapSpringBeans.xml</param>
\end{lstlisting}
This will turn the LDAP Integration on and configure the connection.
\end{enumerate}

\subsubsection*{1.2 Configure System Parameters}
\emph{These system parameters are so vital, that the application will fail
in most instances when not configured properly}
\begin{description}
\item [KIM\_TO\_LDAP\_FIELD\_MAPPINGS] Many different lookup for
  entities storied in LDAP can exist and they are not constrained to
  field name consistency. This means the field names can vary. In the
  event that field names do vary, then there cannot be a one-to-one
  mapping between KIM objects and LDAP objects. Instead, we have to
  have redundant mappings. This system parameter allows us to do just
  that.

When creating the System Parameters, use the following 

\begin{lstlisting}[language=sql,basicstyle=\scriptsize,backgroundcolor=\color{ubergray},caption={trnapp-config.xml},frame=single,breaklines=true]
entityId=uid;principalId=uid;principalName=eduPersonPrincipalName;givenName=sn;principals.principalName=eduPersonPrincipalName;persons.principalName=eduPersonPrincipalName;principals.principalId=uid;principals.active=employeeStatus;lastName=sn;firstName=givenName;employmentInformation.employeeStatus=employeeStatus;employmentInformation.employeeId=emplId,facultyId;names.lastName=sn;names.firstName=givenName;employmentInformation.employeeStatusCode=employeeStatus;
\end{lstlisting}

\item [KIM\_TO\_LDAP\_VALUE\_MAPPINGS] map possible values in the
  lookup form the results from an LDAP query. For example, employee
  active/inactive types have Y, N, and Both values. Each value needs
  to be mapped to a result from the LDAP query results.

\begin{lstlisting}[language=sql,basicstyle=\scriptsize,backgroundcolor=\color{ubergray},caption={trnapp-config.xml},frame=single,breaklines=true]
principals.active.Y=T;principals.active.N=!T;
\end{lstlisting}

\end{description}

\subsubsection*{1.3 Make KIMLdapSpringBeans.xml More Configurable}
We want to be able to modify the KIMLdapSpringBeans.xml which is in
the ldap module of the Rice project. We do not want to have to rebuild
it each time we make a change to the Spring beans.

\begin{enumerate}
\item Copy \textbf{KIMLdapSpringBeans.xml} from
  \textbf{src/main/resources/org/kuali/rice/ldap/config} in the
  \textbf{ldap} module in \textbf{rice-src} to your
  \textbf{src/main/resources} in the \textbf{trnapp} project.

%%\includegraphics[width=\textwidth]{images/Screenshot19.png}

%%\includegraphics[width=\textwidth]{images/Screenshot20.png}

\item Edit the new \textbf{KIMLdapSpringBeans.xml} to look like:

\begin{lstlisting}[numbers=left,language=xml,basicstyle=\scriptsize,backgroundcolor=\color{ubergray},caption={trnapp/src/main/resources/KIMLdapSpringBeans.xml},frame=single,breaklines=true]
	<bean id="kimConstants" class="org.kuali.rice.kim.util.ConstantsImpl">
  	<property name="kimLdapIdProperty"     	value="${rice.ldap.param.kimLdapIdProperty}" />
  	<property name="kimLdapNameProperty"   	value="${rice.ldap.param.kimLdapNameProperty}" />
  	<property name="snLdapProperty"        	value="${rice.ldap.param.snLdapProperty}" />
  	<property name="givenNameLdapProperty" 	value="${rice.ldap.param.givenNameLdapProperty}" />
  	<property name="entityIdKimProperty"   	value="${rice.ldap.param.entityIdKimProperty}" />
  	<property name="employeeMailLdapProperty"  value="${rice.ldap.param.employeeMailLdapProperty}" />
  	<property name="employeePhoneLdapProperty" value="${rice.ldap.param.employeePhoneLdapProperty}" />
  	<property name="defaultCountryCode"    	value="${rice.ldap.param.defaultCountryCode}" />
  	<property name="mappedParameterName"   	value="$mappedParameterName}" />
  	<property name="mappedValuesName"      	value="${rice.ldap.param.mappedValuesName}" />
  	<property name="parameterNamespaceCode"	value="${rice.ldap.param.parameterNamespaceCode}" />
  	<property name="parameterDetailTypeCode"   value="${rice.ldap.param.parameterDetailTypeCode}" />
  	<property name="personEntityTypeCode"  	value="${rice.ldap.param.personEntityTypeCode}" />
  	<property name="employeeIdProperty"    	value="${rice.ldap.param.employeeIdPropertyd}" />
  	<property name="departmentLdapProperty"	value="${rice.ldap.param.departmentNumber}" />  
  	<property name="employeeTypeProperty"  	value="${rice.ldap.param.employeeTypeProperty}" />
  	<property name="employeeStatusProperty"	value="${rice.ldap.param.employeeStatusProperty}" />
  	<property name="affiliationLdapProperty"   value="${rice.ldap.param.affiliationLdapProperty}" />
  	<property name="primaryAffiliationLdapProperty"   value="${rice.ldap.param.primaryAffiliationLdapProperty}" />
  	<property name="defaultCampusCode"     	value="${rice.ldap.param.defaultCampusCode}" />
  	<property name="defaultChartCode"      	value="${rice.ldap.param.defaultChartCode}" />
  	<property name="taxExternalIdTypeCode" 	value="${rice.ldap.param.taxExternalIdTypeCode}" />
  	<property name="externalIdProperty"    	value="${rice.ldap.param.externalIdProperty.externalId}" />
  	<property name="externalIdTypeProperty"	value="${rice.ldap.param.externalIdTypeProperty.externalIdentifierTypeCode}" />
  	<property name="affiliationMappings"
    value="${rice.ldap.param.affiliationMappings}"
      />
  	<property name="employeeAffiliationCodes"  value="${rice.ldap.param.employeeAffiliationCodes}" />
	</bean>
\end{lstlisting}
Notice that I prefixed the parameters with
\textbf{rice.ldap.param}. This is the standard way to create a context
for your parameters. Otherwise, it is unpolished and can be vulnerable
to a name collision.
\item Create a new file called \textbf{ldap-config.xml} in your
  \textbf{trnapp} project. The file will be located in
  \textbf{src/main/resources/META-INF}
  \item Copy the contents of \textbf{KIMLdapSpringBeans.xml} into it.
  \item Then edit the file using search/replace until it looks like a
    standard XML Config. Here is what I had:

\begin{lstlisting}[numbers=left,language=xml,basicstyle=\scriptsize,backgroundcolor=\color{ubergray},caption={trnapp/src/main/resources/KIMLdapSpringBeans.xml},frame=single,breaklines=true]
<config>
    <param name="rice.ldap.param.kimLdapIdProperty"     	>uid</param>
  	<param name="rice.ldap.param.kimLdapNameProperty"   	>eduPersonPrincipalName</param>
  	<param name="rice.ldap.param.snLdapProperty"        	>sn</param>
  	<param name="rice.ldap.param.givenNameLdapProperty" 	>givenName</param>
  	<param name="rice.ldap.param.entityIdKimProperty"   	>entityId</param>
  	<param name="rice.ldap.param.employeeMailLdapProperty"  >mail</param>
  	<param name="rice.ldap.param.employeePhoneLdapProperty" >employeePhone</param>
  	<param name="rice.ldap.param.defaultCountryCode"    	>1</param>
  	<param name="rice.ldap.param.mappedParameterName"   	>KIM_TO_LDAP_FIELD_MAPPINGS</param>
  	<param name="rice.ldap.param.mappedValuesName"      	>KIM_TO_LDAP_VALUE_MAPPINGS</param>
  	<param name="rice.ldap.param.parameterNamespaceCode"	>KR-SYS</param>
  	<param name="rice.ldap.param.parameterDetailTypeCode"   >All</param>
  	<param name="rice.ldap.param.personEntityTypeCode"  	>PERSON</param>
  	<param name="rice.ldap.param.employeeIdProperty"    	>emplId</param>
  	<param name="rice.ldap.param.departmentLdapProperty"	>departmentNumber</param>  
  	<param name="rice.ldap.param.employeeTypeProperty"  	>employeeType</param>
  	<param name="rice.ldap.param.employeeStatusProperty"	>employeeStatus</param>
  	<param name="rice.ldap.param.affiliationLdapProperty"   >affiliationProperty</param>
  	<param name="rice.ldap.param.primaryAffiliationLdapProperty"   >eduPersonPrimaryAffiliation</param>
  	<param name="rice.ldap.param.defaultCampusCode"     	>MC</param>
  	<param name="rice.ldap.param.defaultChartCode"      	>UA</param>
  	<param name="rice.ldap.param.taxExternalIdTypeCode" 	>TAX</param>
  	<param name="rice.ldap.param.externalIdProperty"    	>externalIdentifiers.externalId</param>
  	<param name="rice.ldap.param.externalIdTypeProperty"	>externalIdentifiers.externalIdentifierTypeCode</param>
  	<param name="rice.ldap.param.affiliationMappings"
    >staff=STAFF,faculty=FCLTY,employee=STAFF,student=STDNT,affilate=AFLT</param>
  	<param name="rice.ldap.param.employeeAffiliationCodes"  >STAFF,FCLTY</param>
</config>
\end{lstlisting}
\item Add the \textbf{ldap-config.xml} to your \textbf{/home/rice/kuali/main/dev/trnapp-config.xml}

\begin{lstlisting}[numbers=left,language=xml,basicstyle=\scriptsize,backgroundcolor=\color{ubergray},caption={trnapp-config.xml},frame=single,breaklines=true]
<param
name="config.location">classpath:META-INF/ldap-config.xml</param>
\end{lstlisting}
\item Since we are here, change the
  \textbf{rice.additionalSpringFiles} to refer to our new \textbf{KIMLdapSpringBeans.xml}

\begin{lstlisting}[numbers=left,language=xml,basicstyle=\scriptsize,backgroundcolor=\color{ubergray},caption={trnapp-config.xml},frame=single,breaklines=true]
   	<param
    name="rice.additionalSpringFiles">KIMLdapSpringBeans.xml</param>
\end{lstlisting}
\item Now when we can modify our LDAP configuration from our
  \textbf{trnapp-config.xml}, and not worry about having to
  redistribute the rice-ldap.jar each time.
\end{enumerate}

\addcontentsline{toc}{subsection}{Integration}
\subsection*{2 Integration}
We have already configured our Rice LDAP Integration and implemented
it, but what haven't done is actually integrate it into our Rice and
distribute that yet. As it stands, we probably cannot start our
application because of this missing dependency. We will now do the
integration and set that up.

\subsubsection*{2.1 Add the ldap module to rice}
\begin{enumerate}
\item Open the \textbf{rice-src/pom.xml}.
\item Locate the \textbf{modules} section
\item Add the \verb|<module>ldap</module>|

\begin{lstlisting}[numbers=left,language=xml,basicstyle=\scriptsize,backgroundcolor=\color{ubergray},caption={rice-src/pom.xml},frame=single,breaklines=true]
	<modules>
		<module>api</module>
		<module>impl</module>
		<module>ldap</module>
		<module>web</module>
		<module>sampleapp</module>
		<module>ksb</module>
		<module>kcb</module>
		<module>kns</module>
		<module>kim</module>
		<module>kew</module>
		<module>ken</module>
	</modules>
\end{lstlisting}
\item Now install the module. Execute the following from
  \textbf{rice-src}
\begin{lstlisting}[language=bash,basicstyle=\scriptsize,backgroundcolor=\color{ubergray},caption={Install
  the LDAP Integration Module},frame=single,breaklines=true]
mvn clean install
\end{lstlisting}
\end{enumerate}

\subsubsection*{2.2 Go}
Now you can just start your application. Now you have two tomcat
contexts you must start. One for your Rice application, and one for
CAS. 
\begin{enumerate}
\item Go to Run->Start CAS
\item Go to Run->Start Tomcat
\item Browse to http://localhost:8081/cas/login
\item You are now using LDAP Integration
\end{enumerate}

\addcontentsline{toc}{subsection}{Use Case for Adding Caching}
\subsection*{3 Use Case for Adding Caching}
Numerous LDAP Queries (even in a connection pooled environment) can be
taxing on the system. Caching can solve this. There are different
scenarios for solving the torrent of entity information with a
cache. You can implement your own internal cache within your LDAP
implementation, or you can utilize the builtin caching for KIM.

\newpage
{\setlength{\baselineskip}%
  {0.0\baselineskip}
  \section*{Notes}
  \hrulefill \par}
\addcontentsline{toc}{part}{Exercise 2: Database Setup}
\part*{Exercise 2: \\
Database Setup}

\addcontentsline{toc}{section}{Exercise 2: Database Setup}
{\setlength{\baselineskip}%
  {0.0\baselineskip}
  \section*{\flushright Exercise 2: \\
Database Setup}
  \hrulefill \par}

\addcontentsline{toc}{subsection}{Description}
\subsection*{Description}

\addcontentsline{toc}{subsection}{Goals}
\subsection*{Goals}

\addcontentsline{toc}{subsection}{Instructions}
\subsection*{Instructions}

\newpage
{\setlength{\baselineskip}%
  {0.0\baselineskip}
  \section*{Notes}
  \hrulefill \par}
\addcontentsline{toc}{part}{Exercise 3: Create Business Object Table}
\part*{Exercise 3: \\
Create Business Object Table}

\addcontentsline{toc}{section}{Exercise 3: Create Business Object Table}
{\setlength{\baselineskip}%
  {0.0\baselineskip}
  \section*{\flushright Exercise 3: \\
Create Business Object Table}
  \hrulefill \par}

\addcontentsline{toc}{subsection}{Description}
\subsection*{Description}

\addcontentsline{toc}{subsection}{Goals}
\subsection*{Goals}

\addcontentsline{toc}{subsection}{Instructions}
\subsection*{Instructions}
\subsubsection*{1 Run mysql}
\begin{minted}{sh}
  mysql -u kuldev -p kuldev
\end{minted}
\newpage
{\setlength{\baselineskip}%
  {0.0\baselineskip}
  \section*{Notes}
  \hrulefill \par}
\addcontentsline{toc}{part}{Exercise 4: Create Maintenance Document}
\part*{Exercise 4: \\
Setup Business Object}

\addcontentsline{toc}{section}{Exercise 5: Create Maintenance Document}
{\setlength{\baselineskip}%
  {0.0\baselineskip}
  \section*{\flushright Exercise r: \\
Setup Business Object}
  \hrulefill \par}

\addcontentsline{toc}{subsection}{Description}
\subsection*{Description}

\addcontentsline{toc}{subsection}{Goals}
\subsection*{Goals}

\addcontentsline{toc}{subsection}{Instructions}
\subsection*{Instructions}

\subsubsection*{1 Create Business Object Class}
Create a file in \textbf{work/src/org/kuali/kfs/sys/businessobject/}
called \textbf{FerpaCertification}
\begin{minted}{java}
/*
 * Copyright 2005-2008 The Kuali Foundation
 * 
 * Licensed under the Educational Community License, Version 2.0 (the "License");
 * you may not use this file except in compliance with the License.
 * You may obtain a copy of the License at
 * 
 * http://www.opensource.org/licenses/ecl2.php
 * 
 * Unless required by applicable law or agreed to in writing, software
 * distributed under the License is distributed on an "AS IS" BASIS,
 * WITHOUT WARRANTIES OR CONDITIONS OF ANY KIND, either express or implied.
 * See the License for the specific language governing permissions and
 * limitations under the License.
 */
package org.kuali.kfs.sys.businessobject;


/**
 * Ferpa Certification Business Object
 */
public class FerpaCertification extends PersistableBusinessObjectBase implements Inactivateable {
  
  private Long id;
  private String principalId;
  private String Person;

  /**
  * Default no-arg constructor.
  */
  
  public FerpaCertification() {
        super();
    }

    /**
     * Gets the id attribute.
     * 
     * @return Returns the id.
     */
    public String getId() {
        return id;
    }


    /**
     * Sets the id attribute value.
     * 
     * @param id The id to set.
     */
    public void setId(Long id) {
        this.id = id;
    }


    /**
     * Gets the principalId attribute.
     * 
     * @return Returns the principalId.
     */
    public Long getPrincipalId() {
        return principalId;
    }


    /**
     * Sets the principalId attribute value.
     * 
     * @param principalId The principalId to set.
     */
    public void setPrincipalId(String principalId) {
        this.principalId = principalId;
    }


    /**
     * Gets the person attribute.
     * 
     * @return Returns the person.
     */
    public String getPerson() {
        return person;
    }


    /**
     * Sets the person attribute value.
     * 
     * @param person The person to set.
     */
    public void setPerson(String person) {
        this.person = person;
    }
}
\end{minted}

\subsubsection*{2 Create Business DataDictionary Metadata}
\begin{minted}{xml}
\end{minted}

\subsubsection*{3 Create O/R Mapping}

\begin{minted}{xml}
\end{minted}

\subsubsection*{Add a link to the portal}

\newpage
{\setlength{\baselineskip}%
  {0.0\baselineskip}
  \section*{Notes}
  \hrulefill \par}
\addcontentsline{toc}{part}{Exercise 5: Create Maintenance Document}
\part*{Exercise 5: \\
Create Maintenance Document}

\addcontentsline{toc}{section}{Exercise 5: Create Maintenance Document}
{\setlength{\baselineskip}%
  {0.0\baselineskip}
  \section*{\flushright Exercise 5: \\
Create Maintenance Document}
  \hrulefill \par}

\addcontentsline{toc}{subsection}{Description}
\subsection*{Description}

\addcontentsline{toc}{subsection}{Goals}
\subsection*{Goals}

\addcontentsline{toc}{subsection}{Instructions}
\subsection*{Instructions}

\newpage
{\setlength{\baselineskip}%
  {0.0\baselineskip}
  \section*{Notes}
  \hrulefill \par}
\addcontentsline{toc}{part}{Exercise 6: Basic Routing}
\part*{Exercise 6: \\
Basic Routing}

\addcontentsline{toc}{section}{Exercise 6: Basic Routing}
{\setlength{\baselineskip}%
  {0.0\baselineskip}
  \section*{\flushright Exercise 6: \\
Basic Routing}
  \hrulefill \par}

\addcontentsline{toc}{subsection}{Description}
\subsection*{Description}

\addcontentsline{toc}{subsection}{Goals}
\subsection*{Goals}

\addcontentsline{toc}{subsection}{Instructions}
\subsection*{Instructions}
\begin{minted}{xml}
<data xmlns="ns:workflow" xmlns:fo="http://www.w3.org/1999/XSL/Format" xmlns:xsi="http://www.w3.org/2001/XMLSchema-instance" xsi:schemaLocation="ns:workflow resource:WorkflowData">
  <documentTypes xmlns="ns:workflow/DocumentType" xsi:schemaLocation="ns:workflow/DocumentType resource:DocumentType">
    <documentType>
      <name>FERPA</name>
      <label>Ferpa Certification</label>
      <parent>FSSM</parent>
    </documentType>
  </documentTypes>
</data>
\end{minted}
\newpage
{\setlength{\baselineskip}%
  {0.0\baselineskip}
  \section*{Notes}
  \hrulefill \par}
\addcontentsline{toc}{part}{Exercise 7: Split Node}
\part*{Exercise 7: \\
Split Node}

\addcontentsline{toc}{section}{Exercise 7: Split Node}
{\setlength{\baselineskip}%
  {0.0\baselineskip}
  \section*{\flushright Exercise 7: \\
Split Node}
  \hrulefill \par}

\addcontentsline{toc}{subsection}{Description}
\subsection*{Description}

\addcontentsline{toc}{subsection}{Goals}
\subsection*{Goals}

\addcontentsline{toc}{subsection}{Instructions}
\subsection*{Instructions}

\subsubsection*{1 Create FerpaCertificationMaintainable.java}
\begin{minted}{java}
/*
 * Copyright 2009 The Kuali Foundation
 * 
 * Licensed under the Educational Community License, Version 2.0 (the "License");
 * you may not use this file except in compliance with the License.
 * You may obtain a copy of the License at
 * 
 * http://www.opensource.org/licenses/ecl2.php
 * 
 * Unless required by applicable law or agreed to in writing, software
 * distributed under the License is distributed on an "AS IS" BASIS,
 * WITHOUT WARRANTIES OR CONDITIONS OF ANY KIND, either express or implied.
 * See the License for the specific language governing permissions and
 * limitations under the License.
 */
package org.kuali.kfs.module.ar.document;

import org.apache.commons.lang.StringUtils;
import org.kuali.kfs.module.ar.businessobject.InvoiceRecurrence;
import org.kuali.kfs.sys.KFSConstants;
import org.kuali.kfs.sys.context.SpringContext;
import org.kuali.kfs.sys.document.FinancialSystemMaintainable;
import org.kuali.kfs.sys.document.FinancialSystemMaintenanceDocument;
import org.kuali.rice.kew.exception.WorkflowException;
import org.kuali.rice.kim.bo.Person;
import org.kuali.rice.kim.bo.entity.KimPrincipal;
import org.kuali.rice.kim.service.IdentityManagementService;
import org.kuali.rice.kns.service.DocumentService;
import org.kuali.rice.kns.workflow.service.KualiWorkflowDocument;

public class InvoiceRecurrenceMaintainable extends FinancialSystemMaintainable {

\end{minted}

\subsubsection*{2 Implement answerSplitNodeQuestion in FerpaCertificationMaintainable}

\begin{minted}{java}
@Override
protected boolean answerSplitNodeQuestion(String nodeName) throws UnsupportedOperationException {
}    
\end{minted}

\newpage
{\setlength{\baselineskip}%
  {0.0\baselineskip}
  \section*{Notes}
  \hrulefill \par}

\end{document}
